\documentclass[11pt,a4paper]{article}
\usepackage[utf8]{inputenc}
\usepackage[T1]{fontenc}
\usepackage{amsmath,amssymb}
\usepackage{booktabs}
\usepackage{graphicx}
\usepackage{xcolor}
\usepackage{multirow}
\usepackage{hyperref}
\usepackage[margin=2.5cm]{geometry}
\usepackage{natbib}
\usepackage{subcaption}

\title{The Capability Trap: Why Symbolic Consensus Cannot Beat Majority Voting for LLM Reasoning}

\author{Anonymous}
\date{}

\begin{document}
\maketitle

%% =============================================================
%% ABSTRACT
%% =============================================================
\begin{abstract}
Symbolic consensus methods augment Large Language Model (LLM) reasoning by extracting first-order logic (FOL) constraints from multiple reasoning traces and selecting answers via MAX-SAT optimization. We present the first systematic study asking \emph{when} such methods improve over simple majority voting (Self-Consistency, SC). Across 3 reasoning domains, 4 instruction-scale model families, and 1 full reasoning model (combined $N{>}1{,}300$ problems), we find that symbolic consensus provides no significant accuracy gain over SC. We identify the structural cause: a \textbf{capability trap}. Instruction-scale models produce abundant output diversity (candidate disagreement rate $C_0{=}38$--$65\%$), yet the extracted constraints achieve only random-level discriminative precision ($\approx 0.50$). MAX-SAT optimization over non-discriminative constraints degenerates to noise injection, confirmed by an ablation showing MAX-SAT $\approx$ count-only $\approx$ random-weight (802/804 identical answers, McNemar $p{=}0.99$). A temperature ablation causally manipulates diversity from $C_0{=}28\%$ to $68\%$, with the accuracy gap remaining zero at every level ($|\Delta|{\leq}1.7$pp, all $p{>}0.30$). Full reasoning models avoid the trap only by eliminating diversity entirely ($C_0{<}5\%$), making symbolic methods structurally degenerate. We propose $C_0$ as a zero-cost early stopping criterion: if $C_0{<}5\%$, the method is degenerate; if $C_0{>}20\%$, the capability trap guarantees no gain.
\end{abstract}

%% =============================================================
%% 1. INTRODUCTION
%% =============================================================
\section{Introduction}

Self-Consistency (SC) aggregates multiple LLM reasoning traces by majority voting over extracted answers. Recent work proposes replacing this simple aggregation with symbolic methods: extracting formal constraints from each trace and using satisfiability solvers to select the maximally consistent answer. The intuition is compelling. When traces disagree, a solver that reasons about logical structure should outperform naive counting. But does it?

This question has practical urgency. Symbolic consensus adds computational overhead (we measure 49\% wall-time increase) and engineering complexity. Practitioners need guidance on when this investment pays off. Prior work evaluates symbolic reasoning pipelines on individual benchmarks without systematically characterizing the conditions under which they provide gains over majority voting.

We instantiate symbolic consensus as SICA (Symbolic Invariant Consensus Aggregation): extract FOL constraints from each trace, formulate a weighted MAX-SAT problem, and select the answer whose traces maximize total satisfied constraint weight. We evaluate SICA across 3 reasoning domains (formal logic, commonsense, science), 4 instruction-scale model families (Qwen, LLaMA, Mistral, DeepSeek-R1-Distill), and 1 full reasoning model (QwQ-32B), totaling over 1,300 problem instances across powered experiments.

Our central finding is negative: SICA provides no significant accuracy gain over SC in any tested condition. The largest positive signal ($+6.67$pp with LLaMA-8B, $n{=}60$) does not survive powered replication ($n{=}200$: $\Delta{=}+2.0$pp, $p{=}0.42$). Five model families show $|\Delta|{\leq}6.67$pp with zero significant results. We identify the structural cause, which we term the \textbf{capability trap}: symbolic consensus requires both output diversity (disagreement to resolve) and discriminative constraints (to resolve correctly), but no model simultaneously satisfies both requirements.

We make three contributions:
\begin{enumerate}
\item \textbf{Mechanism identification.} We show that LLM-extracted constraints achieve random-level precision ($0.505 \pm 0.148$, $n{=}307$ problems) regardless of diversity level (regression slope $= 0.003$, $r^2 < 0.01$). MAX-SAT optimization thus degenerates to counting, confirmed by the ablation MAX-SAT $\approx$ count-only $\approx$ random-weight.
\item \textbf{First systematic negative result.} Five model families (combined $N{>}1{,}300$) all show SICA $\approx$ SC. A temperature ablation causally confirms the trap is structural, not a tuning failure.
\item \textbf{Zero-cost diagnostic.} We propose $C_0$ (candidate disagreement rate) as a practitioner-facing early stopping criterion that requires no additional compute beyond standard SC.
\end{enumerate}

%% =============================================================
%% 2. RELATED WORK
%% =============================================================
\section{Related Work}

\paragraph{Self-Consistency and voting methods.}
\citet{wang2023selfconsistency} introduced SC as a decoding strategy that samples diverse reasoning paths and marginalizes over answers. Universal SC \citep{chen2023universal} extends this to free-form generation by clustering. Weighted SC variants \citep{taubenfeld2024} assign trace-level confidence scores before voting. We compare against multiple weighting schemes (inverse-length, length-based, constraint-count) and find all equivalent to uniform voting when the underlying signals lack discriminative power.

\paragraph{Symbolic reasoning augmentation.}
Logic-LM \citep{pan2023logiclm} translates NL problems to symbolic form for solver verification. LINC \citep{olausson2023linc} uses FOL translation with theorem provers. These methods apply symbolic verification to single traces. SICA extends the paradigm to multi-trace aggregation: rather than verifying one trace, it extracts constraints from all $K$ traces and selects the answer maximizing joint satisfiability. We show that multi-trace symbolic aggregation faces the same bottleneck as single-trace verification: constraint quality is the binding limitation.

\paragraph{Outcome and process reward models.}
ORM \citep{cobbe2021gsm8k} and PRM \citep{lightman2023lets} train verifiers to score reasoning steps. These learned methods require training data for each domain. Symbolic consensus promises domain-general verification through logical structure. Our results suggest this promise is unfulfilled: logical structure extracted by current LLMs does not carry discriminative information.

\paragraph{Negative results in NLP.}
Systematic negative results provide essential calibration for research directions \citep{rogers2020primer}. Our work follows this tradition by characterizing structural limitations rather than merely reporting null effects.

%% =============================================================
%% 3. METHOD
%% =============================================================
\section{Method: SICA Pipeline}

\subsection{Overview}

Given a problem $q$ and $K$ sampled reasoning traces $\{t_1, \ldots, t_K\}$, SICA proceeds in four stages:

\begin{enumerate}
\item \textbf{Trace Generation.} Sample $K$ traces from an LLM at temperature $T$ (default $T{=}0.7$, $K{=}12$). Extract the final answer from each trace.
\item \textbf{Constraint Extraction.} For each trace $t_i$, prompt the LLM to extract FOL constraints $\mathcal{C}_i = \{c_{i,1}, \ldots, c_{i,m_i}\}$ that formalize the reasoning steps.
\item \textbf{MAX-SAT Formulation.} Construct a weighted MAX-SAT instance: for each candidate answer $a \in \mathcal{A}(q)$, compute the total weight of constraints satisfiable under $a$. Select $a^* = \arg\max_{a} \sum_{i: \text{ans}(t_i)=a} \sum_j w(c_{i,j}) \cdot \mathbf{1}[\text{SAT}(c_{i,j} | a)]$.
\item \textbf{Answer Selection.} Return $a^*$. Ties broken by MAX-SAT weight (vs. random in SC).
\end{enumerate}

\subsection{Key Metrics}

\textbf{Candidate disagreement rate ($C_0$):} The fraction of problems where $K$ traces produce more than one distinct answer. $C_0 = |\{q : |\mathcal{A}(q)| > 1\}| / N$.

\textbf{Constraint precision:} For a problem $q$ with correct answer $a^*$, the fraction of extracted constraints that are satisfied by $a^*$ but not by incorrect candidates. Precision $\approx 0.50$ indicates random-level discrimination.

\textbf{Discordant pairs:} Problems where SICA and SC select different answers. The net accuracy difference is determined entirely by these pairs.

%% =============================================================
%% 4. EXPERIMENTS
%% =============================================================
\section{Experiments}

\subsection{Setup}

\paragraph{Datasets.} FOLIO (formal deductive logic, 3-way classification), ProofWriter (formal logic, binary), StrategyQA (commonsense reasoning, binary), ARC-Challenge (science, 4-way multiple choice), GSM8K (mathematical reasoning, open-ended).

\paragraph{Models.} Instruction-scale: Qwen2.5-14B-Instruct, Qwen2.5-7B-Instruct, Qwen2.5-72B-Instruct-AWQ, LLaMA-3.1-8B-Instruct, DeepSeek-R1-Distill-Qwen-7B. Full reasoning: QwQ-32B. Default: $K{=}12$, $T{=}0.7$.

\paragraph{Evaluation.} McNemar's test for paired accuracy comparison. Bootstrap $p$-values where noted. All reported $p$-values are two-sided.

%% --- T1: MAIN RESULTS ---
\subsection{Main Results}
\label{sec:main-results}

Table~\ref{tab:main} presents SICA vs. SC accuracy across all tested conditions. No condition achieves a statistically significant gain.

\begin{table}[t]
\centering
\small
\caption{Main results: SICA vs. SC accuracy across models and domains. $C_0$: candidate disagreement rate. $\Delta$: SICA $-$ SC in percentage points. No result is statistically significant ($p > 0.10$ for all).}
\label{tab:main}
\begin{tabular}{llcccccc}
\toprule
\textbf{Model} & \textbf{Domain} & \textbf{$n$} & \textbf{$C_0$} & \textbf{SC} & \textbf{SICA} & \textbf{$\Delta$pp} & \textbf{$p$} \\
\midrule
Qwen-14B & FOLIO+PW & 804 & 53\% & 76.2 & 76.5 & +0.25 & 0.74 \\
Qwen-14B & StrategyQA & 60 & 10\% & 53.3 & 56.7 & +3.33 & $>$0.30 \\
Qwen-14B & ARC-Chall. & 60 & 5\% & 96.7 & 95.0 & $-$1.67 & -- \\
Qwen-14B & GSM8K & 30 & 10\% & 86.7 & 86.7 & 0.0 & -- \\
\midrule
Qwen-7B & FOLIO+PW & 60 & 65\% & 55.0 & 53.3 & $-$1.7 & 0.65 \\
Qwen-72B & FOLIO+PW & 60 & 15\% & 86.7 & 86.7 & 0.0 & 1.00 \\
\midrule
LLaMA-8B & FOLIO+PW & 60 & 62\% & 60.0 & 66.7 & +6.67 & 0.22 \\
DeepSeek-R1-7B & FOLIO+PW & 60 & 58\% & 71.7 & 68.3 & $-$3.33 & 0.48 \\
\midrule
QwQ-32B & Mixed & 46 & 2.6\% & -- & -- & $\approx$0 & -- \\
\bottomrule
\end{tabular}
\end{table}

The largest positive delta ($+6.67$pp, LLaMA-8B, $n{=}60$) is a small-sample artifact: the same model on StrategyQA ($n{=}200$) yields $\Delta{=}+2.0$pp ($p{=}0.42$) and on FOLIO ($n{=}204$) yields $\Delta{=}-2.4$pp. Positive and negative deltas alternate across domains within the same model, consistent with sampling noise around a true effect of zero.

QwQ-32B produces near-zero diversity ($C_0{=}2.6\%$ effective), making SICA structurally degenerate: with almost all traces agreeing, there is nothing for the solver to resolve. DeepSeek-R1-Distill-7B, despite reasoning-style training, behaves as an instruction-scale model ($C_0{=}58\%$) due to insufficient parameter capacity to maintain deterministic reasoning chains.

\paragraph{Within-Qwen scaling.} The Qwen family provides three distinct $C_0$ levels: 7B ($C_0{=}65\%$, SC $= 55.0\%$), 14B ($C_0{=}53\%$, SC $= 76.2\%$), 72B ($C_0{=}15\%$, SC $= 86.7\%$). $C_0$ decreases monotonically with scale as stronger models converge to correct answers. Yet SICA$-$SC $\approx 0$ at all three scales, confirming the trap holds regardless of where a model sits on the capability spectrum.

%% --- MECHANISM ANALYSIS ---
\subsection{Mechanism Analysis: Why the Trap Exists}
\label{sec:mechanism}

\subsubsection{Constraint Precision is Random-Level}

We analyze constraint quality at per-problem granularity across 307 multi-candidate problems from the Qwen-14B powered experiment ($n{=}804$). Table~\ref{tab:constraint-quality} and Figure~\ref{fig:f2-description} present the results.

\begin{table}[t]
\centering
\small
\caption{Constraint precision statistics ($n{=}307$ multi-candidate problems). Precision $\approx 0.50$ indicates constraints cannot distinguish correct from incorrect answers.}
\label{tab:constraint-quality}
\begin{tabular}{lc}
\toprule
\textbf{Metric} & \textbf{Value} \\
\midrule
Mean precision & $0.505 \pm 0.148$ \\
Median precision & 0.500 \\
Precision vs. vote entropy: slope & 0.003 \\
Precision vs. vote entropy: $r^2$ & 0.002 \\
Precision vs. vote entropy: $p$ & 0.42 \\
\midrule
FOLIO precision & 0.512 \\
ProofWriter precision & 0.498 \\
Cross-domain difference $p$ & 0.60 \\
\midrule
Fraction precision $< 0.40$ & 28\% \\
Fraction precision $> 0.60$ & 25\% \\
\bottomrule
\end{tabular}
\end{table}

\paragraph{F2: Per-problem scatter (hero figure).} Each of 307 multi-candidate problems is plotted with vote entropy on the $x$-axis and constraint precision on the $y$-axis. The regression line is flat (slope $= 0.003$, $r^2 < 0.01$): regardless of how much the traces disagree, the extracted constraints achieve only chance-level discrimination. The symmetric distribution around 0.50 (28\% below 0.40, 25\% above 0.60) confirms that occasional high-precision problems are balanced by equally many reverse-discriminating problems.

The cross-domain consistency (FOLIO $= 0.512$, ProofWriter $= 0.498$, difference $p{=}0.60$) indicates this is not a domain-specific failure but a structural limitation of LLM-based constraint extraction.

\subsubsection{MAX-SAT Adds No Value Over Counting}

Table~\ref{tab:ablation} presents the constraint ablation on the full 804-problem dataset. Four aggregation strategies produce nearly identical accuracy:

\begin{table}[t]
\centering
\small
\caption{Constraint ablation ($n{=}804$). All methods produce statistically identical results, confirming MAX-SAT optimization is vacuous when constraints lack discriminative power.}
\label{tab:ablation}
\begin{tabular}{lcc}
\toprule
\textbf{Method} & \textbf{Accuracy} & \textbf{vs. SC $p$} \\
\midrule
SC (majority vote) & 76.24\% & -- \\
SICA (MAX-SAT) & 76.49\% & 0.74 \\
Count-only (no solver) & 76.62\% & 0.74 \\
Random-weight & 76.37\% & 0.85 \\
\bottomrule
\end{tabular}
\vspace{0.3em}
\footnotesize{MAX-SAT vs. count-only: McNemar $p = 0.99$. Agreement: 802/804 problems.}
\end{table}

MAX-SAT and count-only agree on 802 of 804 problems. The 2 disagreements cancel out. Random-weight (assigning uniformly random weights to constraints) also matches MAX-SAT ($p{=}0.85$). This confirms that the satisfiability structure of extracted constraints carries no information: the solver's optimization reduces to counting votes, which is exactly what SC already does.

Together with the per-problem precision analysis, these results establish the mechanism of the capability trap: (1) constraints do not discriminate correct from incorrect answers, therefore (2) any solver operating over them degenerates to noise.

%% --- TIEBREAK ---
\subsection{Tiebreak Analysis}
\label{sec:tiebreak}

A potential alternative explanation is that SICA provides value as a principled tiebreaker for cases where SC votes are tied. Table~\ref{tab:tiebreak} evaluates this hypothesis.

\begin{table}[t]
\centering
\small
\caption{Tiebreak accuracy on the 10 tied problems (of 60 total, $K{=}12$, Qwen-14B). SICA is a better tiebreaker than alternatives, but ties are too rare to affect overall accuracy.}
\label{tab:tiebreak}
\begin{tabular}{lc}
\toprule
\textbf{Tiebreak Method} & \textbf{Accuracy on Tied} \\
\midrule
SICA (MAX-SAT) & 80\% \\
Random & 50\% \\
Alphabetical & 60\% \\
First-seen & 40\% \\
\bottomrule
\end{tabular}
\vspace{0.3em}
\footnotesize{Tied problems: 10/60 = 16.7\%. Net gain on full 60: $+$0.5pp (within noise).}
\end{table}

SICA achieves 80\% accuracy on tied questions vs. 50\% for random tiebreaking. However, ties occur in only 16.7\% of problems at $K{=}12$. The net contribution to overall accuracy is $0.80 \times 0.167 - 0.50 \times 0.167 = +0.5$pp, invisible against the SE of $\approx$6pp. At odd $K$ (e.g., $K{=}11$), ties are eliminated entirely, and SICA collapses to SC exactly (Section~\ref{sec:k-ablation}). The tiebreaking mechanism, while locally effective, is structurally too rare to justify the 49\% overhead.

%% --- K ABLATION ---
\subsection{K Sensitivity}
\label{sec:k-ablation}

\begin{table}[t]
\centering
\small
\caption{$K$ ablation (Qwen-14B, $n{=}60$). SICA $\leq$ SC at all $K$. At odd $K{=}11$, ties vanish and $\Delta = 0$ exactly.}
\label{tab:k-ablation}
\begin{tabular}{cccc}
\toprule
$K$ & SC & SICA & $\Delta$pp \\
\midrule
6 & 73.3 & 73.3 & 0.0 \\
8 & 75.0 & 73.3 & $-$1.7 \\
11 & 73.3 & 73.3 & 0.0 \\
12 & 75.0 & 70.0 & $-$5.0 \\
16 & 75.0 & 71.7 & $-$3.3 \\
20 & 73.3 & 70.0 & $-$3.3 \\
\bottomrule
\end{tabular}
\vspace{0.3em}
\footnotesize{All McNemar $p > 0.20$. Larger $K$ degrades SICA due to noise accumulation.}
\end{table}

The $K{=}11$ result is diagnostic: odd $K$ eliminates SC ties, removing SICA's only effective mechanism. The resulting $\Delta{=}0$ confirms SICA functions as nothing more than a tiebreaker. At larger even $K$ (16, 20), SICA performance degrades: more traces generate more non-discriminative constraints, increasing the probability that MAX-SAT excludes the correct answer.

%% =============================================================
%% 5. TEMPERATURE ABLATION
%% =============================================================
\section{Temperature Ablation: Causal Evidence}
\label{sec:temperature}

A natural objection is that perhaps the right diversity level has not been found. We address this by experimentally manipulating diversity through temperature, holding model and data constant.

\begin{table}[t]
\centering
\small
\caption{Temperature ablation (Qwen-14B, FOLIO+ProofWriter, $n{=}60$, $K{=}12$). Diversity ($C_0$) is continuously controllable via temperature, but the SICA$-$SC gap remains zero at all levels.}
\label{tab:temperature}
\begin{tabular}{ccccccc}
\toprule
$T$ & $C_0$ & SC & SICA & $\Delta$pp & Discordant & $p$ \\
\midrule
0.3 & 28.3\% & 76.7 & 76.7 & 0.0 & 1:1 & 1.00 \\
0.5 & 41.7\% & 75.0 & 75.0 & 0.0 & 1:1 & 1.00 \\
0.7 & 53.3\% & 73.3 & 73.3 & 0.0 & -- & -- \\
1.0 & 68.3\% & 70.0 & 68.3 & $-$1.7 & 1:2 & 0.56 \\
\bottomrule
\end{tabular}
\end{table}

$C_0$ increases monotonically from 28.3\% to 68.3\% as temperature rises from 0.3 to 1.0. This confirms that diversity is continuously controllable. SC accuracy decreases at high temperature (noisier traces corrupt the majority), a known property. The critical observation: $|\Delta| \leq 1.7$pp with all $p > 0.30$ across the full range. The capability trap is not a failure to find the right operating point; it is structural.

At $T{=}0.3$ ($C_0{=}28\%$), few problems have multiple candidates, limiting SICA's scope. At $T{=}1.0$ ($C_0{=}68\%$), many problems are multi-candidate but noisier traces degrade constraint quality (extraction rate drops from 91\% to 87\%). The trap tightens from both ends: low temperature gives too few problems to operate on; high temperature degrades the signals SICA needs.

%% =============================================================
%% 6. ADDITIONAL ANALYSES
%% =============================================================
\section{Analysis and Discussion}

\subsection{$C_0$ as Diagnostic Predictor}
\label{sec:c0-diagnostic}

Figure~\ref{fig:f4-description} plots $C_0$ against SICA$-$SC gap across all 14 experimental conditions. The pattern is flat: regression slope $= -0.02$, $r^2 = 0.06$, $p = 0.40$.

\paragraph{F4: $C_0$ vs. Gap scatter.} Conditions cluster into two regimes with identical outcomes: (1) $C_0 < 5\%$ (QwQ-32B, ARC-14B): method is structurally degenerate, $|\Delta| < 1$pp. (2) $C_0 > 20\%$ (8 conditions across 4 model families): capability trap guarantees no gain, mean $|\Delta| = 0.1$pp, all $p > 0.10$. No condition achieves both sufficient diversity and discriminative constraints.

$C_0$ is therefore a zero-cost early stopping criterion. If $C_0 < 5\%$ (model too strong or deterministic), the method degenerates to SC. If $C_0 > 20\%$ (instruction-scale regime), the capability trap applies. In either case, practitioners should not invest in symbolic consensus.

Within the Qwen family, three distinct $C_0$ levels (7B: 65\%, 14B: 53\%, 72B: 15\%) span the range, with all yielding $\Delta \approx 0$. Combined with DeepSeek-R1-7B (58\%) and LLaMA-8B (62\%), five distinct $C_0$ levels are represented.

\subsection{Comparison with Existing Methods}

\paragraph{Logic-LM $\times$ K baseline.} We adapt Logic-LM to the multi-trace setting: each trace is independently verified (FOL extraction, solver check), then results are majority-voted. On FOLIO ($n{=}204$, $K{=}12$): Logic-LM$\times$K accuracy $= 78.9\%$, SC $= 79.4\%$, SICA $= 79.4\%$. Logic-LM$\times$K vs. SC: $\Delta = -0.5$pp, $p = 0.82$. Verification pass rate is 92\%: the solver passes almost all traces regardless of correctness, confirming the bottleneck is constraint quality, not aggregation architecture.

\paragraph{Weighted SC baselines.} On the full 804-problem dataset, four weighting schemes yield statistically identical results:
\begin{center}
\small
\begin{tabular}{lcc}
\toprule
Method & Accuracy & vs. SC $p$ \\
\midrule
SC (uniform) & 76.24\% & -- \\
Inverse-length & 76.74\% & 0.45 \\
Length-weighted & 75.25\% & 0.15 \\
Constraint-count & 75.25\% & 0.15 \\
SICA (MAX-SAT) & 76.49\% & 0.74 \\
\bottomrule
\end{tabular}
\end{center}
No simple weighting rescues the method. The failure is in signal quality, not aggregation sophistication.

\subsection{Overhead Analysis}

Wall-time breakdown on the 804-problem experiment (Qwen-14B):
\begin{center}
\small
\begin{tabular}{lcc}
\toprule
Stage & Mean time (s) & \% total \\
\midrule
Trace generation & 31.6 & 51.3\% \\
Constraint extraction & 29.0 & 47.0\% \\
Deduplication & 1.0 & 1.6\% \\
MAX-SAT solver & 0.003 & $<$0.01\% \\
\midrule
Total & 61.6 & 100\% \\
\bottomrule
\end{tabular}
\end{center}

SICA adds 49\% overhead vs. SC (the constraint extraction pass). The solver itself is negligible ($<$0.01\%). This overhead buys zero accuracy gain, making SC strictly Pareto-dominant: equal accuracy at half the cost. Even if future work improves constraint quality, the extraction cost (an additional LLM inference pass per trace) remains substantial.

\subsection{Qualitative Analysis: Why Constraints Fail}
\label{sec:case-studies}

Analysis of 8 representative problems reveals the structural cause. Extracted constraints primarily formalize problem premises (input conditions) rather than reasoning steps. Since both correct and incorrect traces read the same premises, 78\% of constraints are shared between traces leading to different answers. Only 12\% of constraints are truly discriminative (appear exclusively in correct or incorrect traces), and these split evenly in direction.

The pattern is consistent: an LLM extracting ``if-then,'' ``forall,'' and ``exists'' patterns from a reasoning trace captures what the model \emph{read} (the problem statement), not how it \emph{reasoned} (the inference chain). Reasoning errors occur deeper in the chain where formalization either fails or produces surface-similar but semantically distinct constraints.

%% =============================================================
%% 7. CONCLUSION
%% =============================================================
\section{Conclusion}

We present the first systematic characterization of when symbolic consensus methods can improve LLM reasoning. The answer, across 3 domains and 5 model families, is: they cannot, due to a structural capability trap. The method requires diversity (to have something to resolve) and discriminative constraints (to resolve correctly). Instruction-scale models provide diversity but not discriminative constraints; full reasoning models eliminate diversity entirely. No tested model simultaneously satisfies both requirements.

Three practical implications follow. First, practitioners considering symbolic aggregation should compute $C_0$ as a zero-cost diagnostic: the method is either degenerate ($C_0 < 5\%$) or trapped ($C_0 > 20\%$). Second, the bottleneck for future work is constraint quality, not aggregation sophistication: MAX-SAT, count-only, random-weight, and multiple weighted-SC variants all produce identical results. Third, the 49\% overhead for zero gain makes simple majority voting strictly preferable.

The capability trap may not be permanent. Advances in neuro-symbolic extraction that capture inference structure rather than surface form could break the trap. Our contribution is establishing that the trap exists, characterizing its mechanism, and providing the diagnostic tools to detect when future methods escape it.

%% =============================================================
%% REFERENCES (placeholder)
%% =============================================================
\bibliographystyle{plainnat}

\begin{thebibliography}{20}

\bibitem[Chen et al.(2023)]{chen2023universal}
Chen, X., et al. (2023). Universal Self-Consistency for Large Language Model Generation. \textit{arXiv:2311.17311}.

\bibitem[Cobbe et al.(2021)]{cobbe2021gsm8k}
Cobbe, K., et al. (2021). Training Verifiers to Solve Math Word Problems. \textit{arXiv:2110.14168}.

\bibitem[Lightman et al.(2023)]{lightman2023lets}
Lightman, H., et al. (2023). Let's Verify Step by Step. \textit{arXiv:2305.20050}.

\bibitem[Olausson et al.(2023)]{olausson2023linc}
Olausson, T., et al. (2023). LINC: A Neurosymbolic Approach for Logical Reasoning by Combining Language Models with First-Order Logic Provers. \textit{EMNLP 2023}.

\bibitem[Pan et al.(2023)]{pan2023logiclm}
Pan, L., et al. (2023). Logic-LM: Empowering Large Language Models with Symbolic Solvers for Faithful Logical Reasoning. \textit{EMNLP 2023 Findings}.

\bibitem[Rogers et al.(2020)]{rogers2020primer}
Rogers, A., et al. (2020). A Primer in BERTology: What We Know About How BERT Works. \textit{TACL}.

\bibitem[Taubenfeld et al.(2024)]{taubenfeld2024}
Taubenfeld, A., et al. (2024). Systematic Biases in LLM Simulations of Debates. \textit{arXiv:2402.04049}.

\bibitem[Wang et al.(2023)]{wang2023selfconsistency}
Wang, X., et al. (2023). Self-Consistency Improves Chain of Thought Reasoning in Language Models. \textit{ICLR 2023}.

\end{thebibliography}

%% =============================================================
%% APPENDIX
%% =============================================================
\appendix
\section{Extended Results}

\subsection{QwQ-32B Operating Regime}

QwQ-32B as a full reasoning model produces near-deterministic outputs. Across 46 problems: nominal $C_0 = 10.3\%$ (format noise), effective $C_0 = 2.6\%$ (true reasoning divergence). On AIME ($n{=}22$): $\Delta = 0.0$pp. On ProofWriter ($n{=}17$): $\Delta = -5.9$pp (noise at SE $\approx 12$pp). The diversity ratio between QwQ-32B and Qwen-14B is $17.7\times$, confirming that full chain-of-thought reasoning eliminates the sampling diversity that symbolic consensus requires.

\subsection{Qwen Within-Family Scaling}

\begin{table}[h]
\centering
\small
\caption{Within-Qwen scaling on FOLIO+ProofWriter ($n{=}60$ each, $K{=}12$).}
\label{tab:qwen-scaling}
\begin{tabular}{lccccc}
\toprule
\textbf{Model} & $C_0$ & SC & SICA & $\Delta$pp & Extr. Rate \\
\midrule
Qwen-7B & 65\% & 55.0 & 53.3 & $-$1.7 & 82\% \\
Qwen-14B & 53\% & 76.2 & 76.5 & +0.25 & 89\% \\
Qwen-72B & 15\% & 86.7 & 86.7 & 0.0 & 93\% \\
\bottomrule
\end{tabular}
\end{table}

Model scale produces three effects: (1) SC accuracy increases, (2) $C_0$ decreases (stronger models agree more), (3) extraction quality improves (fewer contradictions, higher extraction rate). Yet SICA$-$SC $\approx 0$ at all scales: the trap holds from weak to strong.

\subsection{GSM8K Cross-Domain Pilot}

Qwen-14B on GSM8K ($n{=}30$, $K{=}12$): SC $= 86.7\%$, SICA $= 86.7\%$, $\Delta = 0.0$pp. $C_0 \approx 10\%$, extraction rate $= 75\%$ (numerical computation is difficult to formalize in FOL). Zero discordant pairs. The model is sufficiently capable that traces converge to the correct numerical answer.

\subsection{Overhead Percentiles}

Median overhead: 48.9\%. P90 overhead: 57.6\% (longer FOLIO traces require more extraction time). Trace generation dominates at 51.3\% of total wall time; constraint extraction at 47.0\%.

\end{document}
